% ABOUTME: Introduction: why pipeline a processor, what we built, what we contribute and the paper's layout.
% ABOUTME: Owner: Navodit. Numbers are from results/.
\section{Introduction}
\label{sec:intro}

A single-cycle processor completes every instruction in one clock cycle, so
the clock period has to fit the slowest instruction: fetch, register read,
ALU, data memory and write back, one after another. Pipelining cuts this path
into stages with registers between them. Several instructions are then in
flight at once, and each cycle only has to cover one stage, so the clock can
run faster. The cost is hazards: an instruction may need a result that an
older instruction has not written yet, or the processor may fetch the wrong
instructions after a branch \cite{patterson2021}.

In the ideal case a five-stage pipeline is five times faster. This project
asks how much of that gain survives in a real design on real hardware. We
built three processors for the same RV32I subset \cite{riscvspec}: a
single-cycle baseline, a five-stage pipeline that stalls on every data
hazard, and the same pipeline with forwarding. All three were verified
against the same reference model, synthesised for the same FPGA and run on
the same benchmarks, so any difference between them comes from the
microarchitecture alone.

Our contributions are:
\begin{itemize}
\item a single-cycle CPU and a five-stage pipelined CPU in Verilog, with one
module per stage and per pipeline register;
\item a verification flow that compares every register and memory word with a
Python instruction set simulator, on directed hazard tests, benchmarks and
constrained-random programs (248 checks);
\item cycle accurate performance counters and pipeline diagrams drawn from the
hardware's own trace;
\item synthesis and place and route results on a Lattice ECP5 FPGA, with the
critical path of each design explained stage by stage;
\item a CPI, IPC and execution time comparison that separates the gain from a
shorter clock from the losses to pipeline fill, stalls and branch flushes.
\end{itemize}

Section~\ref{sec:background} reviews related work and the performance model.
Sections~\ref{sec:single} and~\ref{sec:pipeline} describe the two designs,
Section~\ref{sec:verification} how we verified them and
Section~\ref{sec:results} what we measured. Section~\ref{sec:conclusion}
concludes and sets out the Phase II plan.